Banks in Chile measure interest rate risk in the banking book under a regulatory text that treats flows indexed to the Unidad de Fomento (UF), a unit of account indexed to prices, as a currency of its own: its own zero curve, shock sizes and post-shock floor, a charge of 2.5 percent of the net UF position, and a pair rule that offsets half of one leg's gain against the other's loss. We implement the treatment on an open engine as five switches and attribute its difference to a single-currency treatment of the same flows by the exact Shapley decomposition, on a synthetic bank (30 books per variant) over the Banco Central de Chile curves of \valDate{}. Treating the UF flows as pesos gives a delta NII \repRatioLo{} to \repRatioHi{} times smaller and a delta EVE \repEveShortLo{} to \repEveShortHi{} percent smaller on the committed book (reproduced); over 30 books the NII ratios are $\mainRatioNiiLng\pm\mainRatioNiiLngCi$, $\mainRatioNiiSht\pm\mainRatioNiiShtCi$ and $\mainRatioNiiOff\pm\mainRatioNiiOffCi$. The indexation charge supplies $\mainChargeShareLng\pm\mainChargeShareLngCi$ percent of the delta NII. For the delta EVE the own curve and the own sizes nearly cancel, refuting the prediction that the pair rule dominates. A bound placing the pair rule between the sum of the legs and the positive parts held in \hOneChecked{} comparisons. The UF floor, binding on no committed curve, raises the delta NII in \floorsRaise{} of \floorsTotal{} comparisons, refuting that it never does. The treatments classify banks differently only at net UF positions of \hFourLo{} to \hFourHi{} percent of level 1 capital.
